\section{Calculus and Polar Functions} \label{sec:polarcalc}

The previous section defined polar coordinates, leading to polar functions. We investigated plotting these functions and solving a fundamental question about their graphs, namely, where do two polar graphs intersect?

We now turn our attention to answering other questions, whose solutions require the use of calculus. A basis for much of what is done in this section is the ability to turn a polar function $r=f(\theta)$ into a set of parametric equations. Using the identities $x=r\cos \theta$ and $y=r\sin \theta$, we can create the parametric equations $x=f(\theta)\cos\theta$, $y=f(\theta)\sin\theta$ and apply the concepts of \autoref{sec:par_calc}.

\subsection{Polar Functions and \texorpdfstring{$\dfrac{\dd y}{\dd x}$}{dy/dx}}

We are interested in the lines tangent to a given graph, regardless of whether that graph is produced by rectangular, parametric, or polar equations. In each of these contexts, the slope of the tangent line is $\frac{\dd y}{\dd x}$. Given $r=f(\theta)$, we are generally \emph{not} concerned with $r\,'=\fp(\theta)$; that describes how fast $r$ changes with respect to $\theta$. Instead, we will use $x=f(\theta)\cos\theta$, $y=f(\theta)\sin\theta$ to compute $\frac{\dd y}{\dd x}$. 

Using \autoref{idea:dydxpar} we have
\[\frac{\dd y}{\dd x} = \frac{\dd y}{\dd\theta}\Big/\frac{\dd x}{\dd\theta}.\]
Each of the two derivatives on the right hand side of the equality requires the use of the Product Rule. We state the important result as a Key Idea.

\begin{keyidea}[Finding $\frac{\dd y}{\dd x}$ with Polar Functions]\label{idea:dydxpol}%
Let $r=f(\theta)$ be a polar function. With $x=f(\theta)\cos\theta$ and $y=f(\theta)\sin\theta$,
\index{polar!functions!finding $\frac{\dd y}{\dd x}$}\index{tangent line}
\[
 \frac{\dd y}{\dd x}
 =\frac{\frac{\dd y}{\dd\theta}}{\frac{\dd x}{\dd\theta}}
 =\frac{\fp(\theta)\sin\theta+f(\theta)\cos\theta}{\fp(\theta)\cos\theta-f(\theta)\sin\theta}.
\]
\end{keyidea}

\youtubeVideo{QTa9OZ4iGPo}{The Slope of Tangent Lines to Polar Curves}

% at some point, someone said `` answer should be ``-2'' not ``-1''. ''  whatever that means
\begin{example}[Finding $\frac{\dd y}{\dd x}$ with polar functions.]\label{ex_polcalc1}%
Consider the limaçon $r=1+2\sin\theta$ on $[0,2\pi]$.
\begin{enumerate}
	\item	Find the rectangular equations of the tangent and normal lines to the graph at $\theta=\pi/4$.
	\item	Find where the graph has vertical and horizontal tangent lines.
\end{enumerate}
\solution
\begin{enumerate}
	\item We start by computing $\frac{\dd y}{\dd x}$. With $\fp(\theta) = 2\cos\theta$, we have
	\begin{align*}
	\frac{\dd y}{\dd x} &= \frac{2\cos\theta\sin\theta + \cos\theta(1+2\sin\theta)}{2\cos^2\theta-\sin\theta(1+2\sin\theta)}\\
	&= \frac{\cos\theta(4\sin\theta+1)}{2(\cos^2\theta-\sin^2\theta)-\sin\theta}.
	\end{align*}
	When $\theta=\pi/4$, $\frac{\dd y}{\dd x}=-2\sqrt{2}-1$ (this requires a bit of simplification). In rectangular coordinates, the point on the graph at $\theta=\pi/4$ is $(1+\sqrt{2}/2,1+\sqrt{2}/2)$. Thus the rectangular equation of the line tangent to the limaçon at $\theta=\pi/4$ is 
\[y=(-2\sqrt{2}-1)\bigl(x-(1+\sqrt{2}/2)\bigr)+1+\sqrt{2}/2 \approx  -3.83 x+8.24.\]
The limaçon and the tangent line are graphed in \autoref{fig:polcalc1}. 
	
	The normal line has the opposite-reciprocal slope as the tangent line, so its equation is 
\[y \approx \frac{1}{3.83}x+1.26.\]
	
	\item		To find the horizontal lines of tangency, we find where $\frac{\dd y}{\dd x}=0$ (when the denominator does not equal 0); thus we find where the numerator of our equation for $\frac{\dd y}{\dd x}$ is 0.
	\[\cos\theta(4\sin\theta+1)=0\quad \Rightarrow \quad \cos\theta=0 \quad \text{or}\quad 4\sin\theta+1=0.\]
	On $[0,2\pi]$, $\cos\theta=0$ when $\theta=\pi/2,\ 3\pi/2$. 

Setting $4\sin\theta+1=0$ gives $\theta=\sin^{-1}(-1/4)\approx -0.2527 = -14.48^\circ$. We want the results in $[0,2\pi]$; we also recognize there are two solutions, one in the 3\textsuperscript{rd} quadrant and one in the 4\textsuperscript{th}. Using reference angles, we have our two solutions as $\theta =3.39$ and $6.03$ radians. The four points we obtained where the limaçon has a horizontal tangent line are given in \autoref{fig:polcalc1} with black-filled dots.\bigskip

To find the vertical lines of tangency, we determine where $\frac{\dd y}{\dd x}$ is undefined by setting the denominator of $\frac{\dd y}{\dd x}=0$ (when the numerator does not equal 0).
\begin{align*}
2(\cos^2\theta -\sin^2\theta)-\sin\theta &= 0 .
\intertext{Convert the $\cos^2\theta$ term to $1-\sin^2\theta$:}
2(1-\sin^2\theta-\sin^2\theta)-\sin\theta &= 0\\
4\sin^2\theta + \sin\theta -2 &= 0.
\intertext{Recognize this as a quadratic in the variable $\sin\theta$. Using the quadratic formula, we have}
\sin\theta &= \frac{-1\pm\sqrt{33}}{8}.
\end{align*}
%
\mtable{The limaçon in \autoref{ex_polcalc1} with its tangent line at $\theta=\pi/4$ and points of vertical and horizontal tangency.}{fig:polcalc1}{\begin{tikzpicture}[alt={Blue limacon with inner loop; red tangent at θ = π/4; filled/open dots show four vertical or horizontal tangent points.}]
\begin{axis}[width=1.16\marginparwidth,tick label style={font=\scriptsize},
axis y line=middle,axis x line=middle,name=myplot,minor x tick num=1,
ymin=-.35,ymax=3.2,xmin=-2.1,xmax=2.1]
\addplot [draw={\colorone},thick, smooth,domain=0:360,samples=60]
 ({cos(x)*(1+2*sin(x))},{sin(x)*(1+2*sin(x))});
\addplot [draw={\colortwo},thick, smooth,domain=0.1:2,samples=2] {-3.82843*x+8.24264};
\filldraw (axis cs: 0,3) circle (1.5pt)
		(axis cs: 0,1) circle (1.5pt)
		(axis cs: -.493,-0.125) circle (1.5pt)
		(axis cs: .483,-0.125) circle (1.5pt);
\filldraw [fill=white,draw=black,thick] (axis cs: 1.76011, 1.31048) circle (1.5pt)
		(axis cs: -1.76011, 1.31048) circle (1.5pt)
		(axis cs: 0.368977, 0.572639) circle (1.5pt)
		(axis cs: -0.369009, 0.578743) circle (1.5pt);
\end{axis}
\node [right] at (myplot.right of origin) {\scriptsize $0$};
\node [above] at (myplot.above origin) {\scriptsize $\pi/2$};
\end{tikzpicture}}%
%
We solve $\sin\theta = \frac{-1+\sqrt{33}}8$ and $\sin\theta = \frac{-1-\sqrt{33}}8$:
\begin{align*}
\sin\theta &=\frac{-1+\sqrt{33}}8 & \sin\theta &= \frac{-1-\sqrt{33}}{8}\\
\theta &= \sin^{-1}\left(\frac{-1+\sqrt{33}}8\right) & \theta &= \sin^{-1}\left(\frac{-1-\sqrt{33}}8\right)\\
\theta &\approx 0.6349 & \theta &\approx -1.0030
\end{align*}
In each of the solutions above, we only get one of the possible two solutions as $\sin^{-1}x$ only returns solutions in $[-\pi/2,\pi/2]$, the 4\textsuperscript{th} and 1\textsuperscript{st} quadrants. Again using reference angles, we have:
\[
\sin\theta = \frac{-1+\sqrt{33}}8 \quad \Rightarrow \quad
\theta \approx 0.6349,\ 2.5067 \text{ radians}
\]
and 
\[
\sin\theta = \frac{-1-\sqrt{33}}8 \quad \Rightarrow \quad
\theta \approx 4.1446,\ 5.2802 \text{ radians.}
\]
These points are also shown in \autoref{fig:polcalc1} with white-filled dots.
\end{enumerate}
\end{example}

When the graph of the polar function $r=f(\theta)$ intersects the pole, it means that $f(\alpha) = 0$ for some angle $\alpha$. Making this substitution in the formula for $\dfrac{\dd y}{\dd x}$ given in \autoref{idea:dydxpol} we see
\[
 \frac{\dd y}{\dd x}
 =\frac{\fp(\alpha)\sin\alpha+f(\alpha)\cos\alpha}{\fp(\alpha)\cos\alpha-f(\alpha)\sin\alpha}
 =\frac{\sin\alpha}{\cos\alpha}
 =\tan\alpha.
\]
%
%When the graph of the polar function $r=f(\theta)$ intersects the pole, it means that $f(\alpha)=0$ for some angle $\alpha$. Thus the formula for $\frac{\dd y}{\dd x}$ in such instances is very simple, reducing simply to \[\frac{\dd y}{\dd x} = \tan \alpha.\]
This equation makes an interesting point. It tells us the slope of the tangent line at the pole is $\tan \alpha$; some of our previous work (see, for instance, \autoref{ex_polar3}) shows us that the line through the pole with slope $\tan \alpha$ has polar equation $\theta=\alpha$. Thus when a polar graph touches the pole at $\theta=\alpha$, the equation of the tangent line at the pole is $\theta=\alpha$.

\begin{example}[Finding tangent lines at the pole]\label{ex_polcalc2}%
Let $r=1+2\sin\theta$, a limaçon. Find the equations of the lines tangent to the graph at the pole.
%
\mtable{Graphing the tangent lines at the pole in \autoref{ex_polcalc2}.}{fig:polcalc2}{\begin{tikzpicture}[alt={Blue looped curve through origin; four symmetric red lines through the pole mark its tangent directions.}]
\begin{axis}[width=1.16\marginparwidth,tick label style={font=\scriptsize},
axis y line=middle,axis x line=middle,name=myplot,
ymin=-.5,ymax=1.37,xmin=-1.1,xmax=1.1]
\addplot [draw={\colorone},thick, smooth,domain=0:360,samples=60]
 ({cos(x)*(1+2*sin(x))},{sin(x)*(1+2*sin(x))});
\addplot [draw={\colortwo},thick, smooth,domain=-1:1,samples=2] {x/sqrt(3)};
\addplot [draw={\colortwo},thick, smooth,domain=-1:1,samples=2] {-x/sqrt(3)};
\end{axis}
\node [right] at (myplot.right of origin) {\scriptsize $0$};
\node [above] at (myplot.above origin) {\scriptsize $\pi/2$};
\end{tikzpicture}}
\solution
We need to know when $r=0$. 
\begin{align*}
1+2\sin\theta &= 0\\
\sin\theta &= -1/2\\
\theta &= \frac{7\pi}{6},\ \frac{11\pi}6.
\end{align*}
Thus the equations of the tangent lines, in polar coordinates, are $\theta = 7\pi/6$ and $\theta = 11\pi/6$. In rectangular form, the tangent lines are $y=\tan(7\pi/6)x=\frac x{\sqrt3}$ and $y=\tan(11\pi/6)x=-\frac x{\sqrt3}$. The full limaçon can be seen in \autoref{fig:polcalc1}; we zoom in on the tangent lines in \autoref{fig:polcalc2}.
\end{example}

\subsection{Area}

When using rectangular coordinates, the equations $x=h$ and $y=k$ defined vertical and horizontal lines, respectively, and combinations of these lines create rectangles (hence the name ``rectangular coordinates''). It is then somewhat natural to use rectangles to approximate area as we did when learning about the definite integral.\index{polar!functions!area}

\mnote{\textbf{Note:} Recall that the area of a sector of a circle with radius $r$ subtended by an angle $\theta$ is $A = \frac12\theta r^2$.
\[
\begin{tikzpicture}[alt={Circular sector with radius r and angle θ; accompanying note states its area A = θ r² /2.},x=30pt,y=30pt,thick]
			\draw (2,0) arc (0:50:2) -- (0,0);
			\draw [] (0,0) -- (2,0) node [pos=.5,below] {$r$};
			\draw [fill=black] (0,0) circle (1pt);
			%\draw (1.95,1.0) node {$s$};
			\draw (0,0) node [shift={(15pt,8pt)}] {$\theta$};
			\end{tikzpicture}
\]}

When using polar coordinates, the equations $\theta=\alpha$ and $r=c$ form lines through the origin and circles centered at the origin, respectively, and combinations of these curves form sectors of circles. It is then somewhat natural to calculate the area of regions defined by polar functions by first approximating with sectors of circles. 

Consider \autoref{fig:polararea} (a) where a region defined by $r=f(\theta)$ on $[\alpha,\beta]$ is given. (Note how the ``sides'' of the region are the lines $\theta=\alpha$ and $\theta=\beta$, whereas in rectangular coordinates the ``sides'' of regions were often the vertical lines $x=a$ and $x=b$.)

Partition the interval $[\alpha,\beta]$ into $n$ equally spaced subintervals as $\alpha=\theta_0<\theta_1<\dotso<\theta_n=\beta$. The radian length of each subinterval is $\Delta\theta = (\beta-\alpha)/n$, representing a small change in angle. The area of the region defined by the $i^{\text{th}}$ subinterval $[\theta_{i-1},\theta_i]$ can be approximated with a sector of a circle with radius $f(c_i)$, for some $c_i$ in $[\theta_{i-1},\theta_i]$. The area of this sector is $\frac12[f(c_i)]^2\Delta\theta$. This is shown in part (b) of the figure, where $[\alpha,\beta]$ has been divided into 4 subintervals. We approximate the area of the whole region by summing the areas of all sectors:
%
\mtable{Computing the area of a polar region.}{fig:polararea}{%
\begin{tikzpicture}[alt={Blue r = f(θ) with dashed spokes; shaded wedge from θ = α to β.}]
\begin{axis}[width=1.16\marginparwidth,tick label style={font=\scriptsize},
axis y line=middle,axis x line=middle,name=myplot,
ymin=-.1,ymax=1.1,xmin=-.1,xmax=1.1,axis equal]
\addplot
 [draw={\colorone},fill={\coloronefill},area style, smooth,domain=18:72,samples=30]
% (axis cs:x:{1+.05*cos(9*x)})
 ({cos(x)*(1+.05*cos(9*x))},{sin(x)*(1+.05*cos(9*x))})
  -- (axis cs:0,0) -- cycle;
\addplot [draw={\colorone},thick, smooth,domain=0:90,samples=30]
 ({cos(x)*(1+.05*cos(9*x))},{sin(x)*(1+.05*cos(9*x))});
\draw [thick,draw={\colorone},] (axis cs:0,0) -- (axis cs: 0.905831, 0.294322)
 node [pos=.7,below,rotate=18,black] {\scriptsize $\theta=\alpha$};
\draw [thick,draw={\colorone},] (axis cs:0,0) -- (axis cs:0.313792, 0.965751)
 node [pos=.7,above,rotate=72,black] {\scriptsize $\theta=\beta$};
\draw [thick,draw={\colorone},dashed]
 (axis cs:0,0) -- (axis cs: 0.862592, 0.528597)
  (axis cs:0,0) -- (axis cs: 0.732107, 0.732107)
   (axis cs:0,0) -- (axis cs: 0.497095, 0.811186);
\draw (axis cs:.8,.85) node {\scriptsize $r=f(\theta)$};
\end{axis}
\node [right] at (myplot.right of origin) {\scriptsize $0$};
\node [above] at (myplot.above origin) {\scriptsize $\pi/2$};
\end{tikzpicture}
\\(a)\\
\begin{tikzpicture}[alt={Red r = f(θ) with dashed spokes; shaded wedge from θ = α to β}]
\begin{axis}[width=1.16\marginparwidth,tick label style={font=\scriptsize},
axis y line=middle,axis x line=middle,name=myplot,
ymin=-.1,ymax=1.1,xmin=-.1,xmax=1.1,axis equal]
\addplot
 [draw={\coloronefill},fill={\coloronefill},area style, smooth,domain=18:72,samples=30]
 ({cos(x)*(1+.05*cos(9*x))},{sin(x)*(1+.05*cos(9*x))}) -- (axis cs:0,0) -- cycle;
\draw [thick,draw={\colorone},] (axis cs:0,0) -- (axis cs:0.313792, 0.965751)
 node [pos=.7,above,rotate=72,black] {\scriptsize $\theta=\beta$};
\addplot [draw={\colortwo},fill={\colortwofill},thick, smooth,domain=18:31.5,samples=30]
 ({.96*cos(x)},{.96*sin(x)}) -- (axis cs:0,0) -- cycle;
\addplot [draw={\colortwo},fill={\colortwofill},thick, smooth,domain=31.5:45,samples=30]
 ({1.05*cos(x)},{1.05*sin(x)}) -- (axis cs:0,0) -- cycle;
\addplot [draw={\colortwo},fill={\colortwofill},thick, smooth,domain=45:58.5,samples=30]
 ({1.0*cos(x)},{1*sin(x)}) -- (axis cs:0,0) -- cycle;
\addplot [draw={\colortwo},fill={\colortwofill},thick, smooth,domain=58.5:72,samples=30]
 ({.96*cos(x)},{.96*sin(x)}) -- (axis cs:0,0) -- cycle;
\draw (axis cs:.8,.85) node {\scriptsize $r=f(\theta)$};
\draw [thick,draw={\colortwo},] (axis cs:0,0) -- (axis cs: 0.905831, 0.294322)
 node [pos=.7,below,rotate=18,black] {\scriptsize $\theta=\alpha$};
\addplot [draw={\colorone},thick, smooth,domain=0:90,samples=30]
 ({cos(x)*(1+.05*cos(9*x))},{sin(x)*(1+.05*cos(9*x))});
\end{axis}
\node [right] at (myplot.right of origin) {\scriptsize $0$};
\node [above] at (myplot.above origin) {\scriptsize $\pi/2$};
\end{tikzpicture}
\\(b)}%
%
\[\text{Area} \approx \sum_{i=1}^n \frac12[f(c_i)]^2\Delta\theta.\]
This is a Riemann sum. By taking the limit of the sum as $n\to\infty$, we find the exact area of the region in the form of a definite integral.

\begin{theorem}[Area of a Polar Region]\label{thm:polar_area}%
Let $f$ be continuous and non-negative on $[\alpha,\beta]$, where $0\leq \beta-\alpha\leq 2\pi$. The area  $A$ of the region bounded by the curve $r=f(\theta)$ and the lines $\theta=\alpha$ and $\theta=\beta$ is 
\[
A=\frac12\int_\alpha^\beta[f(\theta)]^2 \dd\theta
=\frac12\int_\alpha^\beta r^{\,2} \dd\theta
\]
\end{theorem}

The theorem states that $0\leq \beta-\alpha\leq 2\pi$. This ensures that region does not overlap itself, which would give a result that does not correspond directly to the area.

\begin{example}[Area of a polar region]\label{ex_polcalc3}%
Find the area of the circle defined by $r=\cos \theta$. (Recall this circle has radius $1/2$.)
\solution
This is a direct application of \autoref{thm:polar_area}. The circle is traced out on $[0,\pi]$, leading to the integral
%
\mnote{\textbf{Note:} \autoref{ex_polcalc3} requires the use of the integral $\ds\int \cos^2\theta\dd\theta$. This is handled well by using the half angle formula as found in the back of this text. Due to the nature of the area formula, integrating $\cos^2\theta$ and $\sin^2\theta$ is required often. We offer here these indefinite integrals as a time-saving measure.
\begin{align*}
	\int\cos^2\theta\dd\theta &= \frac12\theta+\frac14\sin(2\theta)+C\\
	\int\sin^2\theta\dd\theta &= \frac12\theta-\frac14\sin(2\theta)+C
\end{align*}}
%
\begin{align*}
	\text{Area} &= \frac12\int_0^\pi \cos^2\theta\dd\theta \\
	&= \frac12\int_0^\pi \frac{1+\cos(2\theta)}{2}\dd\theta\\
	&= \frac14\bigl(\theta +\frac12\sin(2\theta)\bigr)\Bigg|_0^\pi\\
	&= \frac\pi4.
\end{align*}
Of course, we already knew the area of a circle with radius $1/2$. We did this example to demonstrate that the area formula is correct.
\end{example}

\begin{example}[Area of a polar region]\label{ex_polcalc4}%
Find the area of the cardioid $r=1+\cos\theta$ bound between $\theta=\pi/6$ and $\theta=\pi/3$, as shown in \autoref{fig:polcalc4}.
\solution
This is again a direct application of \autoref{thm:polar_area}.
%
\mtable{Finding the area of the shaded region of a cardioid in \autoref{ex_polcalc4}.}{fig:polcalc4}{\begin{tikzpicture}[alt={Cardioid sector: blue curve r = f(θ) with shaded slice between θ = 0 and π/6, bounded by two radii.}]
\begin{axis}[width=1.16\marginparwidth,tick label style={font=\scriptsize},
axis y line=middle,axis x line=middle,name=myplot,ytick={1},
ymin=-.7,ymax=1.5,xmin=-.5,xmax=2.1]
\addplot
 [draw={\coloronefill},fill={\coloronefill},area style,smooth,domain=30:60,samples=30]
 ({cos(x)*(1+cos(x))},{sin(x)*(1+cos(x))}) -- (axis cs:0,0) -- cycle;
\addplot [draw={\colorone},thick, smooth,domain=0:360,samples=60]
 ({cos(x)*(1+cos(x))},{sin(x)*(1+cos(x))});
\draw [thick,draw={\colorone}] (axis cs:0,0) -- (axis cs: 1.61603, 0.933013)
 node [pos=.6,below,rotate=30,black] {\scriptsize $\theta=\pi/6$};
\draw [thick,draw={\colorone}] (axis cs:0,0) -- (axis cs:0.75, 1.29904)
 node [pos=.6,above,rotate=60,black] {\scriptsize $\theta=\pi/3$};
\end{axis}
\node [right] at (myplot.right of origin) {\scriptsize $0$};
\node [above] at (myplot.above origin) {\scriptsize $\pi/2$};
\end{tikzpicture}}
%
\begin{align*}
	\text{Area}
	&= \frac12\int_{\pi/6}^{\pi/3} (1+\cos\theta)^2\dd\theta\\
	&= \frac12\int_{\pi/6}^{\pi/3} (1+2\cos\theta+\cos^2\theta)\dd\theta\\
	&= \frac12\left[\theta+2\sin\theta+\frac12\theta
	+\frac14\sin(2\theta)\right]_{\pi/6}^{\pi/3} \\
	&= \frac18\bigl(\pi+4\sqrt{3}-4\bigr).
\end{align*}
\end{example}

\subsubsection{Area Between Curves}

Our study of area in the context of rectangular functions led naturally to finding area bounded between curves. We consider the same in the context of polar functions. \index{polar!functions!area between curves}

\mtable{Illustrating area bound between two polar curves.}{fig:polarea3}{\begin{tikzpicture}[alt={Blue outer r_2 and red inner r_1 enclose shaded region between θ = α and β, inner arc from r_1 to pole.}]
\begin{axis}[width=1.16\marginparwidth,tick label style={font=\scriptsize},
axis y line=middle,axis x line=middle,name=myplot,
ymin=-.1,ymax=1.1,xmin=-.1,xmax=1.1]
\addplot [draw={\colortwo}, thick,smooth,domain=0:90,samples=30]
 ({cos(x)*(.7+.05*sin(6*x))},{sin(x)*(.7+.05*sin(6*x))});
\draw (axis cs:.8,.85) node {\scriptsize $r_2=f_2(\theta)$}
      (axis cs:.2,.8) node {\scriptsize $r_1=f_1(\theta)$};
\addplot [draw={\colorone},thick, smooth,domain=0:90,samples=30]
 ({cos(x)*(1+.05*cos(9*x))},{sin(x)*(1+.05*cos(9*x))});
\addplot[draw={\colorone},fill={\coloronefill},thick, smooth,domain=30:60,samples=30]
 ({cos(x)*(1+.05*cos(9*x))},{sin(x)*(1+.05*cos(9*x))})--(axis cs:0,0)--cycle;
\addplot[draw={\colortwo},fill=white,thick,smooth,domain=30:60,samples=30]
 ({cos(x)*(.7+.05*sin(6*x))},{sin(x)*(.7+.05*sin(6*x))})--(axis cs:0,0)--cycle;
\draw [thick,draw={\colorone},] (axis cs:0,0) -- (axis cs: 0.866025, 0.5)
 node [pos=.4,below,rotate=30,black] {\scriptsize $\theta=\alpha$};
\draw [thick,draw={\colorone},] (axis cs:0,0) -- (axis cs:0.475, 0.822724)
 node [pos=.4,above,rotate=60,black] {\scriptsize $\theta=\beta$};
\end{axis}
\node [right] at (myplot.right of origin) {\scriptsize $0$};
\node [above] at (myplot.above origin) {\scriptsize $\pi/2$};
\end{tikzpicture}}

Consider the shaded region shown in \autoref{fig:polarea3}. We can find the area of this region by computing the area bounded by $r_2=f_2(\theta)$ and subtracting the area bounded by $r_1=f_1(\theta)$ on $[\alpha,\beta]$. Thus
\[
\text{Area}
\ = \ \frac12\int_\alpha^\beta r_2^{\,2}\dd\theta
- \frac12\int_\alpha^\beta r_1^{\,2}\dd\theta
= \frac12\int_\alpha^\beta \bigl(r_2^{\,2}-r_1^{\,2}\bigr)\dd\theta.
\]

\begin{keyidea}[Area Between Polar Curves]\label{idea:area_between_polar}%
The area $A$ of the region bounded by $r_1=f_1(\theta)$ and $r_2=f_2(\theta)$, $\theta=\alpha$ and $\theta=\beta$, where $0\le f_1(\theta)\leq f_2(\theta)$ on $[\alpha,\beta]$, is
\[
A = \frac{1}{2} \int_\alpha^\beta [f_2(\theta)]^2 - [f_1(\theta)]^2\dd\theta
= \frac12\int_\alpha^\beta \bigl(r_2^{\,2}-r_1^{\,2}\bigr)\dd\theta.
\]
\end{keyidea}

\mtable{Finding the area between polar curves in \autoref{ex_polcalc5}.}{fig:polcalc5}{\begin{tikzpicture}[alt={Blue circle r = 1 intersects red cardioid r = 1 − cos θ; right‐hand lens (−π/2 ≤ θ ≤ π/2) is shaded.}]
\begin{axis}[width=1.16\marginparwidth,tick label style={font=\scriptsize},
axis y line=middle,axis x line=middle,name=myplot,axis on top,
ymin=-1.6,ymax=1.6,xmin=-.6,xmax=3.2]
\draw [fill={\coloronefill},draw={\colorone},thick, smooth,domain=0:180,samples=60]  (axis cs:1.5,0)circle(1.5);
% fill=white to cover up what we don't want
\addplot [fill=white,draw={\colortwo}, thick,smooth,domain=0:360,samples=60]
 ({cos(x)*(1+cos(x))},{sin(x)*(1+cos(x))});
% redraw the circle for the part that was covered
\draw [draw={\colorone},thick, smooth,domain=0:180,samples=60]
 (axis cs:1.5,0)circle(1.5);
\end{axis}
\node [right] at (myplot.right of origin) {\scriptsize $0$};
\node [above] at (myplot.above origin) {\scriptsize $\pi/2$};
\end{tikzpicture}}

\begin{example}[Area between polar curves]\label{ex_polcalc5}%
Find the area bounded between the curves $r=1+\cos \theta$ and $r=3\cos\theta$, as shown in \autoref{fig:polcalc5}.
\solution
We need to find the points of intersection between these two functions. Setting them equal to each other, we find:
\begin{align*}
1+\cos\theta &= 3\cos \theta \\
 \cos\theta &=1/2\\
\theta &= \pm \pi/3
\end{align*}
Thus we integrate $\frac12\bigl((3\cos\theta)^2-(1+\cos\theta)^2\bigr)$ on $[-\pi/3,\pi/3]$.
\begin{align*}
	\text{Area}
	&= \frac12\int_{-\pi/3}^{\pi/3} \bigl((3\cos\theta)^2-(1+\cos\theta)^2\bigr)\dd\theta\\
	&= \frac12\int_{-\pi/3}^{\pi/3} \bigl( 8\cos^2\theta-2\cos\theta-1\bigr)\dd\theta \\
	&= \frac12\bigl(2\sin(2\theta) - 2\sin\theta+3\theta\bigr)\Bigg|_{-\pi/3}^{\pi/3}\\
	&= \pi.
\end{align*}
Amazingly enough, the area between these curves has a ``nice'' value.
\end{example}

\begin{example}[Area defined by polar curves]\label{ex_polcalc6}%
Find the area bounded between the polar curves $r=1$ and one petal of $r=2\cos(2\theta)$ where $y>0$, as shown in \autoref{fig:polcalc6}(a).
\solution
%
\mtable{Graphing the region bounded by the functions in \autoref{ex_polcalc6}.}{fig:polcalc6}{%
\begin{tikzpicture}[alt={Blue rose r = sin 4 θ overlaps red three‐leaf r = sin 2 θ; intersection rays at θ = ± π/6, ± π/2.}]
\begin{axis}[width=\marginparwidth,tick label style={font=\scriptsize},
axis y line=middle,axis x line=middle,name=myplot,axis on top,
ymin=-1.15,ymax=1.15,xmin=-.5,xmax=2.2]
\addplot[draw={\coloronefill},fill={\coloronefill},area style,domain=30:45]
 ({cos(x)*(2*cos(2*x))},{sin(x)*(2*cos(2*x))})--(axis cs:{sqrt(3)/2},0)--cycle;
\addplot[draw={\coloronefill},fill={\coloronefill},area style,domain=0:30]
 ({cos(x)},{sin(x)})--(axis cs:0,0)--cycle;
\addplot [draw={\colorone},thick, smooth,domain=0:360,samples=90]
 ({cos(x)*(2*cos(2*x))},{sin(x)*(2*cos(2*x))});
\draw[draw={\colortwo},thick](axis cs:0,0)circle(1);
\end{axis}
\node [right] at (myplot.right of origin) {\scriptsize $0$};
\node [above] at (myplot.above origin) {\scriptsize $\pi/2$};
\end{tikzpicture}
\\(a)\\[10pt]
\begin{tikzpicture}[alt={Close-up: sector α ≤ θ ≤ β between the two curves is filled to illustrate the integrand region.}]
\begin{axis}[width=\marginparwidth,tick label style={font=\scriptsize},
axis y line=middle,axis x line=middle,name=myplot,axis on top,
ymin=-.1,ymax=1,xmin=-.1,xmax=1.2]
\addplot[draw={\colorone},fill={\coloronefill},area style,domain=30:45]
 ({cos(x)*(2*cos(2*x))},{sin(x)*(2*cos(2*x))})--(axis cs:{sqrt(3)/2},0)--cycle;
\addplot[draw={\colortwo},fill={\colortwofill},area style,domain=0:30]
 ({cos(x)},{sin(x)})--(axis cs:0,0)--cycle;
\addplot [draw={\colorone},thick, smooth,domain=0:360,samples=90]
 ({cos(x)*(2*cos(2*x))},{sin(x)*(2*cos(2*x))});
\draw[draw={\colortwo},thick](axis cs:0,0)circle(1);
\draw [thick,dashed] (axis cs: 0,0) -- (axis cs:.866,.5);
\end{axis}
\node [right] at (myplot.right of origin) {\scriptsize $0$};
\node [above] at (myplot.above origin) {\scriptsize $\pi/2$};
\end{tikzpicture}
\\(b)}
%
We need to find the point of intersection between the two curves. Setting the two functions equal to each other, we have
\[
2\cos(2\theta) = 1 \quad \Rightarrow \quad \cos(2\theta) = \frac12
\quad \Rightarrow \quad 2\theta = \pi/3\quad \Rightarrow \quad \theta=\pi/6.
\]
In part (b) of the figure, we zoom in on the region and note that it is not really bounded \emph{between} two polar curves, but rather \emph{by} two polar curves, along with $\theta=0$. The dashed line breaks the region into its component parts. Below the dashed line, the region is defined by $r=1$, $\theta=0$ and $\theta = \pi/6$. (Note: the dashed line lies on the line $\theta=\pi/6$.) Above the dashed line the region is bounded by $r=2\cos(2\theta)$ and $\theta =\pi/6$. Since we have two separate regions, we find the area using two separate integrals.

Call the area below the dashed line $A_1$ and the area above the dashed line $A_2$. They are determined by the following integrals:
\[
A_1 = \frac12\int_0^{\pi/6} (1)^2\dd\theta\qquad
A_2 = \frac12\int_{\pi/6}^{\pi/4} \bigl(2\cos(2\theta)\bigr)^2\dd\theta.
\]
(The upper bound of the integral computing $A_2$ is $\pi/4$ as $r=2\cos(2\theta)$ is at the pole when $\theta=\pi/4$.)

We omit the integration details and let the reader verify that $A_1 = \pi/12$ and $A_2 = \pi/12-\sqrt{3}/8$; the total area is $A = \pi/6-\sqrt{3}/8$.
\end{example}

\subsection{Arc Length}

As we have already considered the arc length of curves defined by rectangular and parametric equations, we now consider it in the context of polar equations. Recall that the arc length $L$ of the graph defined by the parametric equations $x=f(t)$, $y=g(t)$ on $[a,b]$ is
\index{arc length}\index{polar!function!arc length}
\begin{equation}
L = \int_a^b \sqrt{[\fp(t)]^2 + [g\primeskip'(t)]^2}\dd t = \int_a^b \sqrt{[x\primeskip'(t)]^2+[y\primeskip'(t)]^2}\dd t.
\label{eq:polar_arclength}
\end{equation}

Now consider the polar function $r=f(\theta)$. We again use the identities $x=f(\theta)\cos\theta$ and $y=f(\theta)\sin\theta$ to create parametric equations based on the polar function. We compute $x\primeskip'(\theta)$ and $y\primeskip'(\theta)$ as done before when computing $\frac{\dd y}{\dd x}$, then apply \autoeqref{eq:polar_arclength}.

The expression $[x\primeskip'(\theta)]^2+[y\primeskip'(\theta)]^2$ can be simplified a great deal; we leave this as an exercise and state that
\[
[x\primeskip'(\theta)]^2+[y\primeskip'(\theta)]^2 = [\fp(\theta)]^2+[f(\theta)]^2.
\]
This leads us to the  arc length formula.

\begin{keyidea}[Arc Length of Polar Curves]\label{idea:polar_arclength}%
Let  $r=f(\theta)$ be a polar function with $\fp$ continuous on an open interval $I$ containing $[\alpha,\beta]$, on which the graph traces itself only once. The arc length $L$ of the graph on $[\alpha,\beta]$ is
\[L = \int_\alpha^\beta \sqrt{[\fp(\theta)]^2+[f(\theta)]^2}\dd\theta = \int_\alpha^\beta\sqrt{(r\,')^2+ r^2}\dd\theta.\]
\end{keyidea}

\begin{example}[Arc Length of Polar Curves]\label{ex_cardioid_length}%
Find the arc length of the cardioid $r=1+\cos \theta$.
\solution
With $r=1+\cos \theta$, we have $r' = -\sin \theta$. The cardioid is traced out once on $[0,2\pi]$, giving us our bounds of integration. Applying \autoref{idea:polar_arclength} we have
\begin{align*}
	L
	& = \int_0^{2\pi} \sqrt{(-\sin \theta)^2  + (1 + \cos \theta)^2}\dd\theta \\
	& = \int_0^{2\pi} \sqrt{\sin^2 \theta + (1 + 2\cos \theta + \cos^2 \theta)}\dd\theta\\
	& = \int_0^{2\pi} \sqrt{2 + 2\cos \theta}\dd\theta\\
	& = \int_0^{2\pi} \sqrt{2 + 2\cos \theta}\ \frac{\sqrt {2-2\cos \theta}}{\sqrt {2-2\cos \theta}}\dd\theta\\
	& = \int_0^{2\pi} \frac{\sqrt{4 - 4\cos^2 \theta}}{\sqrt {2-2\cos \theta}}\dd\theta\\
	& = 2\int_0^{2\pi} \frac{\sqrt{1 - \cos^2 \theta}}{\sqrt {2-2\cos \theta}}\dd\theta\\
	& = 2\int_0^{2\pi} \frac{\abs{\sin\theta}}{\sqrt {2-2\cos \theta}}\dd\theta
\end{align*}
Since the $\sin \theta > 0$ on $[0, \pi]$ and $\sin \theta < 0$ on $[\pi, 2\pi]$ we separate the integral into two parts
\[
2\int_0^{\pi} \frac{\sin \theta}{\sqrt {2-2\cos \theta}}\dd\theta - 2\int_{\pi}^{2\pi} \frac{\sin \theta}{\sqrt {2-2\cos \theta}}\dd\theta
\]
Using the symmetry of the cardioid and $u$-substitution ($u = 2 - 2\cos \theta$) we simplify the integration to
\begin{align*}
L
& = 4\int_0^{\pi} \frac{\sin \theta}{\sqrt {2-2\cos \theta}}\dd\theta\\
& = 2\int_0^4 \frac{1}{\sqrt u}\dd u\\
& = 4  u^{1/2}\biggr|_0^4 = 8.
\end{align*}
\end{example}

\begin{example}[Arc length of a limaçon]\label{ex_polcalc7}%
Find the arc length of the limaçon $r=1+2\sin\theta$.
\solution
With $r=1+2\sin\theta$, we have $r\,' = 2\cos\theta$. The limaçon is traced out once on $[0,2\pi]$, giving us our bounds of integration. Applying \autoref{idea:polar_arclength}, we have
\begin{align*}
	L
	&= \int_0^{2\pi} \sqrt{(2\cos\theta)^2+(1+2\sin\theta)^2}\dd\theta \\
	&=	\int_0^{2\pi} \sqrt{4\cos^2\theta+4\sin^2\theta +4\sin\theta+1}\dd\theta\\
	&=	\int_0^{2\pi} \sqrt{4\sin\theta+5}\dd\theta\\
	&\approx 13.3649.
\end{align*}
%
\mtable{The limaçon in \autoref{ex_polcalc7} whose arc length is measured.}{fig:polcalc7}{\begin{tikzpicture}[alt={Limacon r = 1 + 2 cos θ plotted; small inner loop (θ around π) and outer loop (whole 0 ≤ θ < 2 π) }]
\begin{axis}[width=1.16\marginparwidth,tick label style={font=\scriptsize},
axis y line=middle,axis x line=middle,name=myplot,axis on top,
ymin=-.5,ymax=3.2,xmin=-2.22,xmax=2.22]
\addplot [draw={\colorone},thick, smooth,domain=0:360,samples=90]
 ({cos(x)*(1+2*sin(x))},{sin(x)*(1+2*sin(x))});
\end{axis}
\node [right] at (myplot.right of origin) {\scriptsize $0$};
\node [above] at (myplot.above origin) {\scriptsize $\pi/2$};
\end{tikzpicture}}
%
The final integral cannot be solved in terms of elementary functions, so we resorted to a numerical approximation. (Simpson's Rule, with $n=4$, approximates the value with $13.0608$. Using $n=22$ gives the value above, which is accurate to 4 places after the decimal.)
\end{example}

\subsection{Surface Area}

The formula for arc length leads us to a formula for surface area. The following Key Idea is based on \autoref{idea:surface_area_parametric}.

\begin{keyidea}[Surface Area of a Solid of Revolution]\label{idea:surface_area_polar}%
Consider the graph of the polar equation $r=f(\theta)$, where $\fp$ is continuous on an open interval containing $[\alpha,\beta]$ on which the graph does not cross itself.
\index{surface area!solid of revolution}\index{polar!function!surface area}\index{integration!surface area}
\begin{enumerate}
	\item The surface area of the solid formed by revolving the graph about the initial ray ($\theta=0$) is:
	\[\text{Surface Area} = 2\pi\int_\alpha^\beta f(\theta)\sin\theta\sqrt{[\fp(\theta)]^2+[f(\theta)]^2}\dd\theta.\]
	\item The surface area of the solid formed by revolving the graph about the line $\theta=\pi/2$ is:
	\[\text{Surface Area} = 2\pi\int_\alpha^\beta f(\theta)\cos\theta\sqrt{[\fp(\theta)]^2+[f(\theta)]^2}\dd\theta.\]
\end{enumerate}
\end{keyidea}

\begin{example}[Surface area determined by a polar curve]\label{ex_polcalc8}%
Find the surface area formed by revolving one petal of the rose curve $r=\cos(2\theta)$ about its central axis (see \autoref{fig:polcalc8}).
\solution
\mtable{Finding the surface area of a rose-curve petal that is revolved around its central axis.}{fig:polcalc8}{%
\begin{tikzpicture}[alt={Red and blue roses r = cos 2 θ and r = sin 2 θ share four petals, alternating every π/4.}]
\begin{axis}[width=1.16\marginparwidth,tick label style={font=\scriptsize},
axis y line=middle,axis x line=middle,name=myplot,axis on top,xtick={-1,1},
ytick={-1,1},ymin=-1.1,ymax=1.1,xmin=-1.3,xmax=1.3]
\addplot [draw={\colorone},thick, smooth,domain=0:45,samples=30]
 ({cos(x)*(cos(2*x))},{sin(x)*(cos(2*x))});
\addplot [draw={\colortwo},thick, smooth,domain=45:360,samples=60]
 ({cos(x)*(cos(2*x))},{sin(x)*(cos(2*x))});
\end{axis}
\node [right] at (myplot.right of origin) {\scriptsize $0$};
\node [above] at (myplot.above origin) {\scriptsize $\pi/2$};
\end{tikzpicture}
\\(a)\\[10pt]
\myincludeasythree{
3Droll=0,
3Dortho=0.009,
3Dc2c=0.41893383860588074 -0.887881875038147 0.1901581883430481,
3Dcoo=65.85308837890625 6.341495513916016 17.98039436340332,
3Droo=85}{\marginparwidth}{One blue petal (0 ≤ θ ≤ pi/2) revolved about its axis, forming a translucent spindle-shaped surface for surface-area computation.}{figures/figpolcalc8a_3D}
\\(b)}
We choose, as implied by the figure, to revolve the portion of the curve that lies on $[0,\pi/4]$ about the initial ray. Using \autoref{idea:surface_area_polar} and the fact that $\fp(\theta) = -2\sin(2\theta)$, we have
\begin{align*}
\text{Surface Area}
&= 2\pi\int_0^{\pi/4} \cos(2\theta)\sin(\theta)
\sqrt{\bigl(-2\sin(2\theta)\bigr)^2+\bigl(\cos(2\theta)\bigr)^2}\dd\theta \\
&\approx 1.36707.
\end{align*}
The integral is another that cannot be evaluated in terms of elementary functions. Simpson's Rule, with $n=4$, approximates the value at $1.36751$.%; with $n=10$, the value is accurate to 4 decimal places.
\end{example}


This chapter has been about curves in the plane. While there is great mathematics to be discovered in the two dimensions of a plane, we live in a three dimensional world and hence we should also look to do mathematics in 3D --- that is, in \emph{space}. The next chapter begins our exploration into space by introducing the topic of \emph{vectors}, which are incredibly useful and powerful mathematical objects.

\printexercises{exercises/09-05-exercises}
